\section{命题概述}
\ChallengeCompany{}提出的命题面向嵌入式软硬件与智能计算的融合，要求连接处理器、操作系统、数据采集、通信和模型推理，形成从芯片集成、驱动开发到感控应用的技术体系。\cite{challenge}据此，本方案把工作重点放在嵌入式推理的开发与部署：让开发者能够说明一个算子计算什么、使用何种数据格式、适合哪个计算后端，并检查部署后的结果是否符合原始定义。

端侧模型通常包含矩阵乘、偏置、激活、量化转换与结果整理等不同类型的运算。矩阵加速器适合承担规整的批量计算，向量单元可处理相同规则下的多元素运算，通用处理器则负责控制流程、输入输出和调度。仅连接这些计算部件仍不足以构成可用的推理系统，还需要统一算子语义、数据布局、缓冲区分配和完成通知，使各阶段能够传递正确的数据。\cite{saturn,gemmini}

本项目采用\ProcessorCore{}、\VectorUnit{}与\AcceleratorName{}组成异构计算结构，以\IntegrationFramework{}管理硬件配置和互连，通过\MemoryController{}访问板上存储。原型载体为现有\BoardModel{}，其 FPGA 型号为\texttt{\FPGAModel}。系统在同一计算结构上分别构建\RealTimeSystem{}与\LinuxDistribution{}运行环境，通过重启和镜像切换验证系统适配，不因操作系统数量增加处理器核心或要求双系统并发运行。\cite{technicalreport}

本方案的主要研究内容是\RuntimeName{}。该体系以明确的算子接口为基础，在电脑端检查模型形状与量化定义，生成后端分配和静态内存计划，在板端依次执行算子并输出带样本标识的结果。新增算子需要同时登记数学定义、支持的数据类型、参考实现和可用后端；未支持的模型内容在转换阶段报告，避免在部署后才发现接口缺失。

工程推进以小型模型的完整推理为起点。先建立独立数值参考，再比较三种后端的单算子和逐层结果，随后开展系统仿真、资源与时序检查、板级启动和两种操作系统下的重复运行。模型输入与结果输出通过适配层连接，后续感知、视觉和通信扩展可复用已验证的推理核心。方案的价值体现在一套能够解释、验证和迁移的开发过程，性能提升与实际应用效果由后续测试确认。
